\usepackage[paperwidth=176mm,paperheight=250mm,inner=21mm,outer=18mm,top=20mm,bottom=20mm,headheight=14pt,headsep=6mm,footskip=10mm]{geometry}
\usepackage{amsmath,amssymb,amsthm,bm,mathtools,booktabs,array,enumitem,float}
\usepackage{tikz,fancyhdr,hyperref}
\usetikzlibrary{arrows.meta,calc,patterns,angles,quotes}
\hypersetup{hidelinks,pdftitle={静电静磁进阶：物理竞赛二轮讲义},pdfauthor={}}
\setlength{\parindent}{2em}\setlength{\parskip}{3pt}\setlength{\emergencystretch}{2em}
\linespread{1.18}\raggedbottom\setcounter{tocdepth}{1}
\ctexset{chapter={name={第,讲},number=\chinese{chapter},format=\LARGE\bfseries\centering,beforeskip=16pt,afterskip=22pt},section={format=\Large\bfseries},subsection={format=\large\bfseries}}
\pagestyle{fancy}\fancyhf{}\fancyhead[LE,RO]{\small\thepage}\fancyhead[LO]{\small 静电静磁进阶}\fancyhead[RE]{\small\leftmark}
\renewcommand{\headrulewidth}{0pt}\fancypagestyle{plain}{\fancyhf{}\fancyfoot[C]{\small\thepage}\renewcommand{\headrulewidth}{0pt}}
\theoremstyle{definition}\newtheorem{example}{例}[chapter]\newtheorem{exercise}{习题}[chapter]
\newcommand{\SolutionStart}{\par\medskip\ifdim\pagetotal>0pt\ifdim\dimexpr\pagegoal-\pagetotal\relax<3\baselineskip\newpage\fi\fi\noindent\textbf{解}\quad\ignorespaces}
\newenvironment{solution}{\SolutionStart}{\par\medskip}
\numberwithin{equation}{chapter}
\newcommand{\dd}{\mathop{}\!\mathrm d}\newcommand{\ee}{\mathrm e}\newcommand{\ii}{\mathrm i}\newcommand{\vct}[1]{\bm{#1}}
\DeclareMathOperator{\arsinh}{arsinh}\DeclareMathOperator{\arcosh}{arcosh}\DeclareMathOperator{\artanh}{artanh}\DeclareMathOperator{\sgn}{sgn}\DeclareMathOperator{\sech}{sech}
\tikzset{axis/.style={->,>=Stealth,line width=.4pt},curve/.style={line width=.7pt},guide/.style={dashed,line width=.35pt},every node/.style={font=\small}}
\setlist[enumerate]{itemsep=3pt,topsep=4pt,leftmargin=2em}
\newcommand{\Needspace}[1]{\par\ifdim\dimexpr\pagegoal-\pagetotal\relax<#1\newpage\fi}
\newcommand{\ke}{k_{\rm e}}
\let\LectureContents\tableofcontents
\renewcommand{\tableofcontents}{\begingroup\small\linespread{1}\selectfont\LectureContents\endgroup}

\providecommand{\ExpansionDiagram}[3]{\begin{figure}[H]\centering\begin{tikzpicture}#3\end{tikzpicture}\caption{#2}\label{#1}\end{figure}}
